\chapter{Thiên Nhiên Dạy Ta Sống (Nature Teaches Us to Live)}
\begin{center}
  \textit{\color{truyengold}``Đi vào rừng một buổi sáng --- bạn học được nhiều hơn cả ngàn trang sách.''}\\[4pt]
  \textit{(One morning in the forest --- you learn more than a thousand pages of books.)}\\[4pt]
  \small\color{truyenblue}--- John Muir ---
\end{center}
\ngancach

\section{Bướm Thoát Kén}
\begin{truyen}{Bướm Thoát Kén}{Bài học từ Thiên Nhiên}
\chuhoa{M}{ột} người đàn ông thấy con bướm đang vật lộn để thoát ra khỏi chiếc kén nhỏ bé; ông thương tình dùng kéo cắt mở kén giúp nó.\\[4pt]
\textit{(A man saw a butterfly struggling to emerge from a tiny cocoon; out of pity he used scissors to cut open the cocoon to help it.)}

Con bướm chui ra dễ dàng --- nhưng thân mình phù nề, đôi cánh nhỏ teo và không bao giờ bay được.\\[4pt]
\textit{(The butterfly emerged easily --- but its body was swollen, its wings tiny and shriveled, and it could never fly.)}

Ông không biết rằng chính cuộc vật lộn qua lỗ nhỏ của chiếc kén là cách thiên nhiên ép chất dinh dưỡng từ thân vào cánh --- sự gian khó chính là điều kiện để cánh trưởng thành.\\[4pt]
\textit{(He did not know that the very struggle through the cocoon's narrow opening is nature's way of forcing fluid from the body into the wings --- the hardship itself is what conditions the wings to mature.)}

\end{truyen}
\begin{baihoc}
  Cố giúp ai đó tránh khỏi đấu tranh cần thiết là vô tình cướp đi sức mạnh họ cần để bay.\\[4pt]
  \textit{(Trying to spare someone a necessary struggle is inadvertently robbing them of the strength they need to fly.)}
\end{baihoc}

\section{Đàn Ngỗng Trời Bay Theo Hình Chữ V}
\begin{truyen}{Đàn Ngỗng Trời Bay Theo Hình Chữ V}{Thiên nhiên --- Khoa học Khí động học}
\chuhoa{K}{hi} ngỗng trời bay theo hình chữ V, mỗi con tạo ra luồng gió nâng cho con bay phía sau; cả đàn bay được xa hơn 71\% so với khi bay một mình.\\[4pt]
\textit{(When geese fly in a V formation, each bird creates an updraft for the one behind; the whole flock flies 71\% farther than any bird could flying alone.)}

Khi con đầu đàn mệt, nó bay lui phía sau; một con khác tự động bay lên dẫn đầu.\\[4pt]
\textit{(When the lead goose tires, it falls back; another automatically moves forward to lead.)}

Con nào bị bệnh hay bị thương, hai con khác rời đàn bay xuống cùng --- ở lại cho đến khi con đó hồi phục hoặc qua đời.\\[4pt]
\textit{(If any goose is sick or wounded, two others leave the flock and fly down to stay --- remaining until that goose recovers or dies.)}

Bài học từ đàn ngỗng: lãnh đạo là phục vụ --- chia sẻ gánh nặng là nghĩa vụ --- không bỏ lại ai phía sau là đạo đức.\\[4pt]
\textit{(The lesson from the geese: leadership is service --- sharing burdens is duty --- leaving no one behind is ethics.)}

\end{truyen}
\begin{baihoc}
  Tập thể đoàn kết luôn đi xa hơn cá nhân tài giỏi nhất --- vì sức mạnh cộng hưởng không thể đo bằng số học.\\[4pt]
  \textit{(A united team always travels further than the most talented individual --- because synergistic strength cannot be calculated arithmetically.)}
\end{baihoc}

\section{Cây Sồi Trăm Năm}
\begin{truyen}{Cây Sồi Trăm Năm}{Thiên nhiên --- Triết lý Phương Tây}
\chuhoa{M}{ột} cây sồi trăm tuổi sống sót qua hàng chục cơn bão không phải vì thân to nhất --- mà vì rễ ăn sâu nhất.\\[4pt]
\textit{(A hundred-year-old oak survives dozens of storms not because it has the largest trunk --- but because it has the deepest roots.)}

Những năm đầu đời, cây sồi con trông có vẻ phát triển chậm chạp; thực ra phần lớn năng lượng của nó đang đổ vào rễ phía dưới --- chuẩn bị cho những thập kỷ đứng vững về sau.\\[4pt]
\textit{(In its early years, a young oak appears to grow slowly; in reality most of its energy is flowing into roots below --- preparing for decades of standing firm ahead.)}

Con người cũng vậy: tuổi trẻ đầu tư vào nền tảng --- tri thức, phẩm cách, rèn luyện --- thì về sau mới đứng vững trước mọi bão tố cuộc đời.\\[4pt]
\textit{(People are the same: youth that invests in foundations --- knowledge, character, discipline --- will stand firm against all of life's storms ahead.)}

\end{truyen}
\begin{baihoc}
  Muốn đứng cao, trước tiên phải ăn rễ sâu --- thành công bền vững được xây từ nền tảng vững chắc.\\[4pt]
  \textit{(To stand tall, first grow deep roots --- lasting success is built on solid foundations.)}
\end{baihoc}

\section{Hạt Sen Ngàn Năm Vẫn Nảy Mầm}
\begin{truyen}{Hạt Sen Ngàn Năm Vẫn Nảy Mầm}{Khoa học --- Trung Quốc, 1995}
\chuhoa{N}{ăm} 1995, các nhà khoa học tìm thấy những hạt sen trong một đầm khô ở Trung Quốc --- được xác định có niên đại hơn 1.000 năm tuổi.\\[4pt]
\textit{(In 1995, scientists found lotus seeds in a dry pond in China --- dated to over 1,000 years old.)}

Họ đặt những hạt sen đó vào nước ấm và chờ.\\[4pt]
\textit{(They placed those seeds in warm water and waited.)}

Hai ngày sau, những hạt sen ngàn năm tuổi bắt đầu nảy mầm.\\[4pt]
\textit{(Two days later, the thousand-year-old seeds began to sprout.)}

Tiềm năng sống không bao giờ chết nếu điều kiện phù hợp chưa đến --- nó chỉ chờ đợi, kiên nhẫn, ngàn năm nếu cần.\\[4pt]
\textit{(The potential for life never dies if the right conditions have not yet arrived --- it simply waits, patiently, a thousand years if necessary.)}

\end{truyen}
\begin{baihoc}
  Đừng vội kết luận rằng một hạt giống đã chết --- có thể nó chỉ đang chờ đúng thời điểm để nảy mầm.\\[4pt]
  \textit{(Do not rush to conclude a seed is dead --- perhaps it is only waiting for the right moment to sprout.)}
\end{baihoc}
